% =====================================================================
%  capitulos/cap02.tex — Capítulo 2. Las reglas de la gramática
%  Autor: Augusto M. Trombetta
%  Fuente: Capitulo_02_normalizado.docx
% =====================================================================
\capitulo{Las reglas de la gramática}{Augusto M. Trombetta}

\section{Introducción}

Las reglas de la gramática no parecen ser todas iguales: unas definen las
características constitutivas propias de los objetos de una lengua, otras
establecen normas de corrección y etiqueta propias del uso social de esa
lengua. La diferencia entre ambos tipos de reglas es correlativa de la
distinción moderna entre gramáticas científicas (descriptivas y explicativas de
los fenómenos lingüísticos observables) y gramáticas tradicionales (normativas
y prescriptivas de los usos lingüísticos correctos). En un caso, las reglas se
descubren, se postulan como conjeturas vinculadas con el funcionamiento interno
de la lengua (por ejemplo, \forma{la oración consta de un sujeto y un
predicado}); en el otro, las reglas se fijan, se establecen como normas
externas relativas a las convenciones sociales de uso de la lengua en un lugar
y época determinados (por ejemplo, \forma{no debe haber coma entre el sujeto y
el predicado}). Este capítulo gira en torno a las reglas de la gramática. En
cuanto a su organización, en primer lugar, se esquematizan las diferentes
concepciones de la gramática y se contrastan los criterios evaluativos de
corrección y gramaticalidad presentes en estas; en segundo lugar, se
especifican las nociones de gramaticalidad y aceptabilidad incorporadas en las
gramáticas científicas, incluyendo las dimensiones normativas como parte de las
condiciones de aceptabilidad; finalmente, se vincula la noción de
gramaticalidad con las variedades lectales que presentan las lenguas.

\section{Concepciones de la gramática}

Como se dijo en el \capref{1}, tradicionalmente, la gramática ha sido
considerada una disciplina sujeta a reglas que rigen el buen uso, es decir, una
disciplina \termino{normativa} o \termino{prescriptiva}. Más modernamente, la
gramática, ya sea considerada en un sentido estático como \textit{la
descripción de un cierto estado de lengua}, ya sea definida en un sentido
dinámico como \textit{el dispositivo mental interno que hace posible articular
las estructuras propias de una lengua}, se entiende como constituida por un
sistema de relaciones, reglas y operaciones, que se deben descubrir para
comprender la conducta lingüística: en este sentido, se trata de una disciplina
\termino{descriptiva} o \termino{científica}.

La valoración de una combinación de la lengua (sílaba, palabra, sintagma,
oración) como \termino{correcta} o \termino{incorrecta} solo adquiere sentido en
el marco de una gramática normativa o prescriptiva, con reglas que controlan
externamente las emisiones lingüísticas. En contraste, una gramática
descriptiva o científica, con reglas que explican internamente el
funcionamiento de la lengua, valora las estructuras de la lengua como
\termino{gramaticales} o \termino{agramaticales}. Esquematizando, se pueden
definir dos maneras complementarias de concebir la gramática y de tipificar las
reglas que forman parte de esta, asociadas a su vez a dos maneras de evaluar
las formas propias del lenguaje humano:

\begin{table}[htbp]
  \centering
  \begin{tabularx}{\linewidth}{@{}>{\itshape}l C C@{}}
    \toprule
    Concepciones de la gramática & Normativa o prescriptiva & Descriptiva o científica \\
    \midrule
    Tipos de reglas          & Regulativas (externas)  & Constitutivas (internas) \\
    Criterios de evaluación  & Corrección/\hspace{0pt}Incorrección & Gramaticalidad/\hspace{0pt}Agramaticalidad \\
    \bottomrule
  \end{tabularx}
  \caption{Concepciones de la gramática}\label{tab:02-concepciones}
\end{table}

\begin{ejercicio}\label{ej:02-definiciones}
Considere las siguientes definiciones 1-5, representativas de las dos
concepciones señaladas sobre la gramática y, posteriormente, complete con los
valores «\mas» o «\menos» (\sig{sí} o \sig{no}) el cuadro presentado más abajo,
como se muestra en el ejemplo resuelto~1.

\begin{exe}
  \ex La \textit{gramática de la actual lengua castellana} no es otra cosa que
      \textit{el conjunto ordenado de las reglas de lenguaje que vemos
      observadas en los escritos o conversación de las personas doctas que
      hablan el castellano o español} \parencite[1]{salva1852}.
  \ex La \textit{gramática} de una lengua es el arte de hablarla correctamente,
      esto es, conforme al buen uso, que es el de la gente educada. Se prefiere
      este uso porque es el más uniforme en las varias provincias y pueblos que
      hablan una misma lengua, y por lo tanto el que hace que más fácil y
      generalmente se entienda lo que se dice; al paso que las palabras y frases
      propias de la gente ignorante varían mucho de unos pueblos y provincias a
      otros, y no son fácilmente entendidas fuera de aquel estrecho recinto en
      que las usa el vulgo \parencite[§§\,1-2]{bello2002}.
  \ex La lingüística estática o descripción de un estado de lengua se puede
      llamar \textit{gramática} en el sentido muy preciso, y por lo demás usual
      [\ldots] donde se trata de un objeto complejo y sistemático que pone en
      juego valores coexistentes. La gramática estudia la lengua como sistema
      de medios de expresión; quien dice gramatical dice sincrónico y
      significativo \parencite[223]{saussure1945}.
  \ex {[La lengua]} Es un tesoro depositado por la práctica del habla en los
      sujetos que pertenecen a una misma comunidad, un sistema gramatical
      virtualmente existente en cada cerebro, o, más exactamente, en los
      cerebros de un conjunto de individuos, pues la lengua no está completa en
      ninguno, no existe perfectamente más que en la masa
      \parencite[57]{saussure1945}.
  \ex La gramática de L será, pues, un ingenio que genere todas las secuencias
      gramaticales de L y ninguna de las agramaticales
      \parencite[27]{chomsky1974a}.
\end{exe}

\begin{tabularx}{\linewidth}{@{}L *{5}{Y{.8cm}}@{}}
  \toprule
  De acuerdo con las definiciones dadas, la gramática de una lengua\ldots & 1 & 2 & 3 & 4 & 5 \\
  \midrule
  \ldots\ remite al uso lingüístico de ciertos grupos sociales.               & \mas   &&&& \\
  \ldots\ representa la descripción de un sistema de valores lingüísticos.    & \menos &&&& \\
  \ldots\ constituye un mecanismo capaz de producir estructuras lingüísticas. & \menos &&&& \\
  \ldots\ implica una regulación modélica, exterior al individuo que la usa.  & \mas   &&&& \\
  \ldots\ supone un sistema o mecanismo internalizado, interior al individuo. & \menos &&&& \\
  \ldots\ conforma una disciplina normativa o prescriptiva.                   & \mas   &&&& \\
  \ldots\ conforma una disciplina descriptiva o científica.                   & \menos &&&& \\
  \bottomrule
\end{tabularx}

\textit{Nota}: La columna correspondiente a 1 muestra el ejemplo resuelto.
\end{ejercicio}

\begin{ejercicio}\label{ej:02-searle}
Considere las siguientes citas de \textcite[55]{searle1964},\footnote{En su
ensayo de 1969, \textit{Actos de habla}, Searle retoma su distinción entre
reglas regulativas y constitutivas, dando a entender que las reglas
gramaticales son regulativas: «Los diferentes lenguajes humanos, en la medida
en que son intertraducibles, pueden considerarse como plasmaciones
convencionales diferentes de las mismas reglas subyacentes. El hecho de que en
francés pueda hacerse una promesa diciendo \enquote*{Je promets} y que en
castellano pueda hacerse diciendo \enquote*{Yo prometo}, es un asunto de
convención. Pero el hecho de que una emisión de un dispositivo de prometer
cuente como (bajo condiciones apropiadas) la asunción de una obligación, es un
asunto de reglas y no un asunto de convenciones del francés o del castellano»
\parencite[48-49]{searle1994}. Es decir, las lenguas naturales aparecen como
convenciones que plasman o materializan las reglas subyacentes que gobiernan
los actos de habla. En el mismo sentido, considerando el dominio estricto de
las estructuras de una lengua, se mantiene la distinción introducida por Searle
entre constitutivo y regulativo: una gramática subyacente, formada por reglas
constitutivas que son la condición de posibilidad de las estructuras
lingüísticas, y una normativa gramatical que regula convencionalmente el uso de
tales estructuras posibles.} donde se establece la diferencia entre
\termino{reglas regulativas} y \termino{reglas constitutivas}. Posteriormente
complete con los valores «\mas» o «\menos» (\sig{sí} o \sig{no}) el cuadro
presentado más abajo, como se muestra en el ejemplo resuelto~(a).

\begin{exe}
  \ex Algunas reglas regulan formas de comportamiento previamente existentes.
      Por ejemplo, las reglas de etiqueta en la mesa regulan el evento de
      comer, pero el evento de comer existe con independencia de estas reglas.
      Existen reglas, por otro lado, que no se limitan a regular, sino que
      crean o definen nuevas formas de comportamiento: las reglas del ajedrez,
      por ejemplo, no se limitan a regular una actividad previamente existente
      llamada jugar al ajedrez; tales reglas, por así decirlo, crean la
      posibilidad o definen esa actividad.
  \ex La distinción que estoy tratando de hacer fue adelantada por la
      distinción de Kant entre principios regulativos y constitutivos, así que
      adoptemos su terminología y describamos nuestra distinción como una
      distinción entre reglas regulativas y constitutivas. Las reglas
      regulativas regulan actividades cuya existencia es independiente de tales
      reglas; las reglas constitutivas constituyen (y también regulan) formas de
      actividad cuya existencia depende lógicamente de estas reglas.
\end{exe}

\begin{tabularx}{\linewidth}{@{}L Y{1.7cm} Y{1.7cm}@{}}
  \toprule
  ¿Qué tipo de regla encierran los siguientes enunciados? & Regulativa & Constitutiva \\
  \midrule
  a) En una oración, todo predicado tiene un sujeto. & \menos & \mas \\
  b) No corresponde decir «pienso de que mañana empiezo la Facultad». & & \\
  c) No se dice «haiga», se dice «haya». & & \\
  d) «Si seguís hablando, la vas a pasar mal» se entiende como amenaza. & & \\
  e) Toda estructura lingüística contiene un núcleo. & & \\
  f) Usar el léxico lunfardo queda prohibido en la radiodifusión argentina. & & \\
  \midrule
  ¿Qué tipo de regla ofrecen las gramáticas normativas o prescriptivas? & & \\
  ¿Qué tipo de regla presentan las gramáticas descriptivas o científicas? & & \\
  ¿Qué tipo de regla es responsable de la regulación interna de una lengua? & & \\
  ¿Qué tipo de regla es responsable de la regulación externa de una lengua? & & \\
  \bottomrule
\end{tabularx}

\textit{Nota}: La fila correspondiente a (a) muestra el ejemplo resuelto.
\end{ejercicio}

\section{Corrección y gramaticalidad}

Según lo presentado hasta acá, las gramáticas normativas o prescriptivas
normalmente se encuentran asociadas a un conjunto de reglas que hacen a los
aspectos regulativos de la lengua, o sea, a las condiciones del buen uso de las
formas lingüísticas. Tales reglas regulativas, en la medida en que representan
regulaciones conscientes de actividades o conductas cuya existencia no depende
de ellas, constituyen factores externos a la estructura de la lengua, que
frecuentemente se justifican de acuerdo con criterios históricos y «que son
producto de las relaciones sociales\ldots» \parencite[29]{bosque2009}. Aunque el
tradicional concepto de \termino{corrección}, aspecto central de las gramáticas
normativas o prescriptivas, fue objeto de reconsideración a lo largo del
siglo~\textsc{xx}, conviene observar que la existencia de «gramáticas
normativas [\ldots] es esperable, e incluso necesaria por razones educativas, y
en general, sociales» \parencite[29]{bosque2009}.

Podemos ilustrar las reglas regulativas con las que figuran en \textit{El buen
uso del español}, una obra normativa elaborada por la Real Academia Española y
la Asociación de Academias de la Lengua \parencite{rae2013}, en la que se
pueden encontrar las siguientes recomendaciones:

\begin{exe}
  \ex\label{ex:02-coma} No se escribe coma entre el sujeto y el predicado,
      independientemente del orden que adopten, ni aunque el sujeto sea
      complejo o una subordinada sustantiva (p.~66).
  \ex\label{ex:02-loismo} Recibe el nombre de loísmo la utilización de las
      formas masculinas de acusativo \forma{lo/los} en lugar de las de dativo
      (\forma{le/les}) [\ldots] Aunque aparece ocasionalmente en el español
      clásico, nunca entró de lleno en la lengua literaria. Hoy se registra
      solamente en algunas zonas de Castilla. El loísmo está fuertemente
      desprestigiado, por lo que se recomienda evitarlo en todos los niveles de
      la lengua (p.~261).
  \ex\label{ex:02-dequeismo} Se llama dequeísmo el uso incorrecto de la
      secuencia \forma{de que} cuando la preposición \forma{de} no está
      justificada [\ldots] El dequeísmo se ha extendido de forma desigual en
      los países hispanohablantes, más en la lengua oral que en la escrita y
      algo más en el español americano que en el europeo, aunque se documenta
      ampliamente en ambos. A pesar de esta difusión, el dequeísmo constituye
      una incorrección gramatical que se recomienda evitar (p.~382).
\end{exe}

Asociado a las gramáticas descriptivas o científicas, el concepto de
\termino{gramaticalidad}, de incorporación más reciente en el campo de la
lingüística, apunta a los aspectos constitutivos de la estructura del lenguaje,
a los factores internos que caracterizan su naturaleza o propiedades. Siguiendo
a \textcite{bosque2009}, las reglas constitutivas, por representar pautas «que
constituyen los fundamentos de algún sistema» (p.~28), ofrecen los «principios
que articulan el sistema lingüístico» (p.~29). Ajeno a la dicotomía normativa
entre formas \termino{correctas} y formas \termino{incorrectas}, el
descubrimiento de tales principios fundamenta tanto la aceptación de las
secuencias bien formadas (\termino{gramaticales}) como la exclusión de las
secuencias mal formadas (\termino{agramaticales}).

Las siguientes observaciones o descripciones extraídas de la
\textit{Introducción a la sintaxis del español}, una obra orientada a la
enseñanza \parencite{camacho2018}, permiten ejemplificar las reglas
constitutivas:

\begin{exe}
  \ex\label{ex:02-haiga} ¿Es lo mismo «gramatical» que «correcto»? No, el
      primero es un concepto cognitivo, el segundo es un concepto social. Por
      ejemplo, en español latinoamericano se oye frecuentemente la forma verbal
      \forma{haiga} (\forma{ojalá que haiga comida en la fiesta}), y algunas
      personas dicen que solo la primera forma es «correcta». Esto refleja una
      jerarquía social: los hablantes que dicen \forma{haya} tienen más
      prestigio social que los que dicen \forma{haiga}. Pero esto no tiene nada
      que ver con la gramaticalidad: la gramática de los hablantes produce
      \forma{haiga} y lo reconoce como posible, por lo tanto, es gramatical
      (p.~17).
  \ex\label{ex:02-faxear} Más arriba hablábamos de los
      morfemas\footnote{El concepto de morfema y los demás términos de la
      morfología se desarrollan en el \capref{4}.} derivativos como morfemas
      que permiten relacionar palabras con una raíz común, pero con
      significados parcialmente distintos, como \forma{mente} y \forma{mental}.
      Este mecanismo también nos sirve para formar palabras nuevas. Por
      ejemplo, a partir de la palabra \forma{fax}, se puede introducir el verbo
      \forma{faxear}, y a partir de \forma{chat}, \forma{chatear} (p.~75).
  \ex\label{ex:02-recursividad} Una de las propiedades más interesantes del
      lenguaje humano es la recursividad,\footnote{Esta noción se presentó en
      el \capref{1}.} que nos permite expresar un número potencialmente
      infinito de pensamientos usando recursos limitados (p.~311).
\end{exe}

\begin{ejercicio}\label{ej:02-correccion}
Considere las siguientes secuencias lingüísticas y evalúe su corrección y
gramaticalidad completando con los valores «\mas» o «\menos» (\sig{sí} o
\sig{no}) el cuadro presentado más abajo, como se muestra en el ejemplo
resuelto~1. Luego responda las preguntas sobre 4a y 4b.

\begin{tabularx}{\linewidth}{@{}L Y{1.5cm} Y{1.5cm}@{}}
  \toprule
  ¿Qué evaluación merecen las siguientes secuencias lingüísticas? & Correcta & Gramatical \\
  \midrule
  1. a. Pierden derechos y carecen por completo de toda protección. & \mas & \mas \\
  \phantom{1.} b. Pierden y carecen por completo. & \menos & \menos \\
  \phantom{1.} c. Pierden derechos y carecen de toda protección. & \mas & \mas \\
  \addlinespace
  2. a. Las diputadas están saliendo y están entrando a otra reunión. & & \\
  \phantom{2.} b. Las diputadas están saliendo de una reunión y están entrando a otra. & & \\
  \addlinespace
  3. a. Trajimos tres remeras iguales para que no haya diferencias de color. & & \\
  \phantom{3.} b. Trajimos tres remeras iguales para que no haiga diferencias de color. & & \\
  \phantom{3.} c. Trajimos tres remeras iguales para que no hay diferencias de color. & & \\
  \addlinespace
  4. a. Tu mujer trabaja todo el día, parece \textit{blecansain}. & & \\
  \phantom{4.} b. Tu \textit{jermu} trabaja todo el día, parece incansable. & & \\
  \bottomrule
\end{tabularx}

\textit{Nota}: El bloque correspondiente a 1 muestra el ejemplo resuelto.

\begin{resuelto}[Aclaraciones sobre el ejemplo resuelto]
1. a. [Pierden [derechos]] y [carecen [por completo] [de toda protección]] /
b. [Pierden Ø] y [carecen [por completo] Ø]. Si bien \forma{pierden} puede
usarse sin COD (\forma{derechos}), si se omite el CR de \forma{carecer}
(\forma{de toda protección}), la oración resulta agramatical. / c. [Pierden
[derechos]] y [carecen Ø [de toda protección]]. La omisión del circunstancial
\forma{por completo}, al no ser un elemento exigido por el verbo, no impide que
la oración siga siendo gramatical. En el contraste de a-b-c se advierte que la
buena formación (gramaticalidad) de las oraciones no depende de la mera
acumulación de palabras: ciertos elementos pueden ser requeridos por la
semántica del verbo (argumentos), frente a otros que resultan facultativos
(adjuntos). Tomado en un sentido amplio, se puede decir que la incorrección de
b depende de su falta de gramaticalidad («es incorrecta porque es
agramatical»), pero no se puede hacer la afirmación opuesta.
\end{resuelto}

\begin{resuelto}[Preguntas sobre 4a y 4b]
¿Es lo mismo \forma{incansable} que \forma{blecansain}? [Esta pregunta la podría
contestar cualquier hablante nativo del español.] ¿Es lo mismo \forma{mujer}
que \forma{jermu}? [Esta pregunta parece apropiada para que la conteste
cualquier hablante nativo del español de la Argentina.]
\end{resuelto}
\end{ejercicio}

\begin{ejercicio}\label{ej:02-registros}
Considere la siguiente cita de \textcite[30]{bosque2009} y conteste las dos
preguntas que se formulan más abajo.

\begin{cita}
Los diferentes modos de comportarse lingüísticamente en una determinada
comunidad se denominan \textsc{registros} o \textsc{estilos}. Los registros se
parecen a los atuendos, los modales, los gestos y otras manifestaciones del
comportamiento social en que son relativos a las situaciones y resultan
apropiados o inapropiados en función de ellas [\ldots] Al igual que sucede con
los atuendos, una parte del conocimiento del idioma afecta a las
\textsc{condiciones de uso} que se asocian con muchas palabras y con algunas
construcciones sintácticas. En muchos casos, los juicios que cabe hacer sobre
las construcciones han de tener en cuenta la variante del sistema gramatical en
la que resultan apropiados.
\end{cita}

\begin{consignas}
  \item ¿Hay algún contraste entre las respuestas esperadas para los puntos 3a
        y 3b del ejercicio~\ref{ej:02-correccion}? [En general, ¿hay algún
        contraste entre las secuencias gramaticales y las secuencias
        correctas?]
  \item Al señalar las condiciones de uso, los registros y los estilos de la
        lengua, ¿los autores están aludiendo al concepto de corrección o al de
        gramaticalidad?
\end{consignas}
\end{ejercicio}

\section{Gramaticalidad y aceptabilidad}\label{sec:02-aceptabilidad}

Los conceptos de \termino{gramaticalidad} y \termino{aceptabilidad} forman
parte del desarrollo de la gramática generativa\footnote{Ver \capref{1}, §\ref{sec:01-generativa} y
§\ref{sec:01-proyeccion}.} y derivan de las obras de Chomsky. Inicialmente, ambas nociones
aparecían muy relacionadas, ya que la idea de \textit{aceptable} venía a
explicar la de \textit{gramaticalidad}, según se puede leer en
\textit{Estructuras sintácticas} \parencite{chomsky1974a}:

\begin{cita}
El propósito fundamental del análisis lingüístico de una lengua L es el de
separar secuencias gramaticales que son oraciones de L, de las secuencias
agramaticales que no son oraciones de L, y estudiar la estructura de las
secuencias gramaticales [\ldots] Una manera de probar (\enquote*{test}) si una
gramática propuesta para L es adecuada, es determinar si las oraciones que
genera son o no son realmente gramaticales, es decir, aceptables para un
hablante nativo (p.~27).
\end{cita}

Son dos conceptos diferentes, como el propio Chomsky se encarga de destacar en
\textit{Aspectos de la teoría de la sintaxis} \parencite{chomsky1970}, donde se
estableció una clara diferencia al relacionarlos con las nociones de
\termino{competencia} (conocimiento lingüístico implícito del hablante de una
lengua) y \termino{actuación} (uso concreto de la lengua por parte de un
hablante). En todo caso, la gramaticalidad es condición necesaria de la
aceptabilidad, pero no es el único factor que interviene en ella:

\begin{cita}
La distinción relacionada entre locuciones \enquote*{aceptables} y
\enquote*{gramaticales} aparece también en SS-LSLT [\textit{Syntactic
structures} y \textit{The logical structure of linguistic theory}], aunque la
terminología es introducida, claramente sólo años después. Podemos tomar
\enquote*{aceptabilidad} como un concepto de la teoría de la actuación, y
\enquote*{gramaticalidad} como una noción relacionada de la teoría de la
competencia. Las oraciones generadas por la gramática directamente son las
oraciones \enquote*{gramaticales} (o \enquote*{gramaticales en el más alto
grado}) de la lengua. Entre ellas hay oraciones que son verdaderas, falsas,
divertidas, ininteligibles, triviales, oraciones que no tienen sentido, etc.
Puede ser muy apropiado, en algunas ocasiones, construir y usar locuciones que
se desvían de las reglas gramaticales en un respecto u otro, y puede muy bien
haber principios generales para interpretar tales oraciones [\ldots] Podemos
decir que la gramática \enquote*{genera derivativamente} las oraciones que se
desvían de las reglas de modos específicos, y que estas oraciones tienen los
grados más bajos de gramaticalidad. Las oraciones son aceptables (o quizá
aceptables bajo circunstancias concretas) si son convenientes, apropiadas,
adecuadas al propósito del momento, etc. La gramática de la competencia
contribuye a determinar la aceptabilidad, pero este concepto implica muchos
otros factores \parencite[6-7]{chomsky1974a}.\pendiente{¿Chomsky 1970
[1965]? El texto atribuye la cita a \textit{Aspectos}}
\end{cita}

Según \textcite[32]{bosque2009}, la \termino{gramaticalidad}, como expresión de
los puros factores gramaticales que «corresponden propiamente al sistema
lingüístico», aparece como una extensión de la teoría de la
\termino{competencia}, el conocimiento tácito (no explícito) del sistema de
reglas constitutivas de su lengua. En tanto, la \termino{aceptabilidad}, que es
la expresión de factores tanto gramaticales como extragramaticales, se presenta
como una extensión de la teoría de la \termino{actuación}, que conjuga factores
de naturaleza diversa (no únicamente gramatical), entre los que se incluyen la
memoria, la consciencia, la rapidez de procesamiento, el cansancio, el estado
de ánimo, el sueño, el nerviosismo, la ingesta de alcohol, etcétera.

\begin{table}[htbp]
  \centering
  \begin{tabularx}{\linewidth}{@{}L L@{}}
    \toprule
    \textbf{Aceptabilidad} & \textbf{Gramaticalidad} \\
    \midrule
    Podemos tomar \enquote*{aceptabilidad} como un concepto de la teoría de la
    \textit{actuación} [\ldots]
      & [\ldots] y \enquote*{gramaticalidad} como una noción relacionada de la
        teoría de la \textit{competencia}. \\
    \addlinespace
    Las oraciones son aceptables [\ldots] si son convenientes, apropiadas,
    adecuadas al propósito del momento, etc. \textit{La gramática de la
    competencia contribuye a determinar la aceptabilidad, pero este concepto
    implica muchos otros factores}.
      & Las oraciones \textit{generadas por la gramática directamente} son las
        oraciones \enquote*{gramaticales} [\ldots] de la lengua. Entre ellas hay
        oraciones que son verdaderas, falsas, divertidas, ininteligibles,
        triviales, oraciones que no tienen sentido, etc. \\
    \bottomrule
  \end{tabularx}
  \caption{Aceptabilidad vs. gramaticalidad}\label{tab:02-aceptabilidad}
\end{table}

\begin{ejercicio}\label{ej:02-aceptabilidad}
Considere las siguientes secuencias lingüísticas y evalúe su gramaticalidad y
aceptabilidad completando el siguiente cuadro con los valores «\mas» o
«\menos» (\sig{sí} o \sig{no}), como se muestra en el ejemplo resuelto~1.

\begin{tabularx}{\linewidth}{@{}L Y{1.5cm} Y{1.5cm}@{}}
  \toprule
  ¿Qué evaluación merecen las siguientes secuencias lingüísticas? & Gramatical & Aceptable \\
  \midrule
  1. a. Estoy seguro de que algunas personas no saben que ciertos periodistas
  afirman que determinados funcionarios no entienden que la razón de su
  nombramiento no viene dada por el currículum que declaran. & \mas & \menos \\
  \phantom{1.} b. Estoy seguro de que algunas personas no lo saben. Ciertos
  periodistas afirman que determinados funcionarios no entienden una cosa
  simple: la razón de su nombramiento no viene dada por el currículum que
  declaran. & \mas & \mas \\
  \addlinespace
  2. La madre del primo de la novia del hermano de mi suegra está contenta. & & \\
  \addlinespace
  3. a. «Toda fuerza armada o reunión de personas que se atribuya los derechos
  del pueblo y peticione a nombre de éste, comete delito de sedición»
  (Constitución Nacional, art.~22). & & \\
  \phantom{3.} b. Toda fuerza armada o reunión de chabones que se atribuya los
  derechos del pueblo y peticione a nombre de éste, comete delito de sedición. & & \\
  \bottomrule
\end{tabularx}

\textit{Nota}: El bloque correspondiente a 1 muestra el ejemplo resuelto.

\begin{resuelto}[Aclaraciones sobre el ejemplo resuelto]
Tanto 1a como 1b son gramaticales. Frente a 1b, la aceptabilidad de 1a se ve
comprometida fuertemente por la cantidad de oraciones subordinadas, incrustadas
unas dentro de las otras, lo que dificulta el procesamiento cognitivo de la
oración final: [Estoy seguro de [que algunas personas no saben [que ciertos
periodistas afirman [que determinados funcionarios no entienden [que la razón
de su nombramiento no viene dada por el currículum [que declaran]]]]]].
\end{resuelto}
\end{ejercicio}

\begin{ejercicio}\label{ej:02-ingrid}
Considere los efectos del alcohol (¿acaso el fernet?) como uno de los factores
no gramaticales que intervienen en la aceptabilidad de las emisiones
lingüísticas. En un video que circuló bastante durante 2015 («El hit del
momento: ¡Dale Ingrid dale!»,
\url{https://www.youtube.com/watch?v=u1AAwMVNxs8}) se puede ver a un joven que
entona una canción cuya letra va cambiando.

\begin{consignas}
  \item ¿Las oraciones de la canción son gramaticales? [Tenga presente la
        estructura de las oraciones: todas las repeticiones que entona.]
  \item ¿Las oraciones de la canción son aceptables? [Tenga presente la
        organización léxica de las oraciones: las primeras dos repeticiones
        frente a las siguientes.]
\end{consignas}
\end{ejercicio}

\section{Gramaticalidad, aceptabilidad y corrección}

El criterio de corrección, o sea la valoración entre \termino{formas correctas}
y \termino{formas incorrectas}, propio de las gramáticas tradicionales de tipo
normativo o prescriptivo vinculadas de manera típica a la enseñanza, se puede
entender como uno de los componentes incluidos en la noción de aceptabilidad (la
dicotomía entre \termino{estructuras aceptables} y \termino{estructuras no
aceptables}) que, según se señaló, es un concepto vinculado con la teoría de la
actuación dentro de un enfoque formal de la gramática. En este sentido,
volviendo sobre la cita de Chomsky (§\ref{sec:02-aceptabilidad}), se puede advertir que la
gramaticalidad (la \textit{buena formación} oracional, es decir, la forma que
adquieren las estructuras lingüísticas según las pautas constitutivas internas
de la gramática) es la condición necesaria para que se materialicen las
estructuras oracionales, con independencia de factores tales como su valor de
verdad y sentido, su elegancia y corrección, su farragosidad o su elocuencia,
etcétera. Queda planteada, en consecuencia, una distinción fundamental entre la
noción de gramaticalidad, asociada con la buena formación de las construcciones
lingüísticas, y la de aceptabilidad, dentro de la que se incluye el criterio de
corrección, y vinculada a los factores sociales e individuales que intervienen
en la actuación lingüística.

\begin{table}[htbp]
  \centering
  \begin{tabularx}{\linewidth}{@{}L L@{}}
    \toprule
    \textbf{Competencia} & \textbf{Actuación} \\
    \midrule
    Reglas y principios constitutivos de la gramática de una lengua natural
      & Factores regulativos y limitantes (adecuación, estilo, corrección,
        memoria, etc.) \\
    \addlinespace
    Gramaticalidad (condición de buena formación de las construcciones de una
    lengua)
      & Aceptabilidad (depende de la situación social, del procesamiento
        mental, etc.) \\
    \bottomrule
  \end{tabularx}
  \caption{Competencia vs. actuación}\label{tab:02-competencia}
\end{table}

Finalmente, la teoría de la competencia (el conocimiento interno del conjunto de
principios y reglas de la gramática de una lengua natural) y la teoría de la
actuación (la puesta en uso por parte del individuo de su competencia
lingüística), con sus conceptos asociados de gramaticalidad y aceptabilidad, no
deja de remitir a la distinción saussureana entre la lengua (el sistema de
valores opositivos internalizado por los miembros de una comunidad lingüística)
y el habla (el uso del sistema de la lengua por parte de un individuo). La
diferencia entre ambas perspectivas pasa por el carácter dinámico que tiene la
gramática interna en Chomsky, que es considerada un dispositivo generativo,
frente al carácter estático y descriptivo que tiene la lengua en Saussure, que
es planteada como un conjunto de relaciones y oposiciones entre los elementos
que la integran.

En lo que sigue se va a contrastar la consistencia formal del concepto de
gramaticalidad con distintas dimensiones que pueden incidir en la aceptabilidad
de una secuencia bien formada. Para ello se van a revisar algunas cuestiones
interpretativas, como la claridad, la verosimilitud o la significatividad, como
así también nociones vinculadas con el uso social del lenguaje, como la
elegancia.

\subsection{Ambigüedad}\label{sec:02-ambiguedad}

El concepto de ambigüedad refiere a secuencias bien formadas que generan más de
una interpretación, como, por ejemplo, \forma{Mi libro}, que puede hacer
referencia a \sig{el libro que poseo} o a \sig{el que escribí}. En las
gramáticas normativas tradicionales se prescribía evitar las expresiones
ambiguas cuando son anfibológicas, es decir, cuando se trata de formas con
múltiples interpretaciones que no resultan claras. Según
\textcite[35]{bosque2009}, que abordan el concepto desde una perspectiva formal,
una secuencia posee ambigüedad gramatical o estructural «cuando puede
corresponderle más de una estructura sintáctica, y por tanto más de una
interpretación». Desde esta perspectiva, el concepto de ambigüedad es
constitutivo de la gramática y apunta a que algunas oraciones por lo menos
admiten «dos análisis sintácticos en el sistema gramatical del español, lo que
da lugar a otros tantos significados». Por ejemplo, una oración como
\exref{ex:02-secarropa} admite dos lecturas, según la manera en que se
interpreten los vínculos entre los elementos que la constituyen:\footnote{La
noción de constituyente se introduce en el \capref{5}.}

\begin{exe}
  \ex\label{ex:02-secarropa} Por favor, cuando termine de usar el secarropa,
      limpie el filtro de pelusas.
\end{exe}

En \exref{ex:02-secarropa}, se puede considerar que [\forma{el filtro}] y
[\forma{de pelusas}] son dos modificadores independientes del verbo
\exref{ex:02-lectura1-a} y que, por lo tanto, sus posiciones son
intercambiables, o bien que [\forma{de pelusas}] modifica a [\forma{el filtro}]
\exref{ex:02-lectura2-a}, por lo que constituyen una unidad. Cada una de estas
lecturas implica un significado diferente de la oración
\exref{ex:02-lectura1-b} y \exref{ex:02-lectura2-b}:

\begin{exe}
  \ex\label{ex:02-lectura1}
    \begin{xlist}
      \ex\label{ex:02-lectura1-a} Por favor, cuando termine de usar el
          secarropa, limpie [el filtro] [de pelusas].
      \ex\label{ex:02-lectura1-b} \textit{Por favor, cuando termine de usar el
          secarropa, limpie [de pelusas] [el filtro].}
    \end{xlist}
  \ex\label{ex:02-lectura2}
    \begin{xlist}
      \ex\label{ex:02-lectura2-a} Por favor, cuando termine de usar el
          secarropa, limpie [el filtro de pelusas].
      \ex\label{ex:02-lectura2-b} \textit{Por favor, cuando termine de usar el
          secarropa, limpie [el filtro [de pelusas]].}
    \end{xlist}
\end{exe}

La ambigüedad gramatical, también conocida como homonimia estructural, puede
darse tanto en el nivel sintáctico como morfológico.\footnote{La ambigüedad u
homonimia estructural en el nivel sintáctico se trata en el \capref{5}, §\ref{sec:05-homonimia} y,
en el nivel morfológico, se estudia en el \capref{4}, §\ref{sec:04-homonimia}.} Pero además de
gramatical, la ambigüedad puede ser léxica, es decir, que puede estar
relacionada con el doble significado que tienen muchas palabras, como se
muestra en el siguiente meme en el que el humor surge a partir de las dos
interpretaciones que la palabra \forma{privado} admite en el mismo contexto:

\begin{exe}
  \ex\label{ex:02-privado} El barrio era privado: privado de carne, de leche, de
      pan\ldots, privado de todo (Diego Maradona, entrevista con J.~Guinzburg,
      \url{https://www.facebook.com/watch/?v=133231901641222}).
    \begin{xlist}
      \ex \forma{privado}\textsubscript{1} adj. que no es propiedad pública o
          estatal, sino que pertenece a particulares.
      \ex \forma{privado}\textsubscript{2} participio pasado del verbo
          \forma{privar}: despojado de (\textit{DLE}, en línea).
    \end{xlist}
\end{exe}

\begin{ejercicio}\label{ej:02-ambiguedad}
Considere las siguientes oraciones y conteste las preguntas que se formulan
abajo, como se muestra en el ejemplo resuelto~1.

\begin{exe}
  \ex Me trajeron esta escultura de Francia.
  \ex Antes de salir a navegar, acomodaron correctamente las velas.
  \ex Yo no recibiría a tu hermano en mi oficina.
\end{exe}

\begin{consignas}
  \item Las tres oraciones anteriores pueden tener más de una interpretación.
        Realice paráfrasis que pongan de relieve los diferentes significados y
        explique de qué tipo de ambigüedad se trata en cada caso.
  \item ¿Observa alguna mala formación de estructura gramatical? ¿Se podría
        resolver de alguna manera el problema de ambigüedad observado en (a)?
\end{consignas}

\begin{resuelto}
1. (a) La oración está bien formada y es gramatical, aunque resulte ambigua, ya
que, fuera de contexto, admite múltiples significados. En este caso, la
ambigüedad se relaciona con factores sintácticos: a. [Me trajeron [esta
escultura [de porcelana [de Francia]]]] / a'. Me trajeron esta escultura
(hecha) de porcelana de Francia. / b. [Me trajeron [esta escultura [de
porcelana]] [de Francia]] / b'. Me trajeron desde Francia esta escultura de
porcelana. (b) La ambigüedad se podría resolver reestructurando el orden de los
constituyentes, como en b'.\pendiente{El ítem 1 no dice «de porcelana»}
\end{resuelto}
\end{ejercicio}

\subsection{Elegancia en la expresión: cortesía/descortesía, redundancia}

La regla de elegancia «no constituye un principio gramatical», ya que al
infringirla no se obtiene una secuencia agramatical, sino «una construcción poco
apropiada o poco elegante» \parencite[36]{bosque2009}. En este sentido, el mayor
o menor grado de elegancia presente en una construcción «se distingue de su
gramaticalidad» y se puede decir que una regla de elegancia constituye «más
bien una norma de cortesía», en tanto se trata de una forma de expresarse que
manifiesta \textit{atención, respeto o afecto} hacia otro (\textit{DLE}, en
línea). Lo contrario de la cortesía, la descortesía, se encuentra cuando alguien
se manifiesta mediante expresiones descalificantes, como las que son frecuentes
en redes sociales:

\begin{exe}
  \ex\label{ex:02-carroza} Ahí lo tenés al pelotudo (meme de amplio uso en redes
      sociales, creado a partir de una escena del filme \textit{Esperando la
      carroza}, de A.~Doria, 1985).
\end{exe}

Asimismo, también se producen expresiones poco elegantes cuando se incurre en
una redundancia o repetición innecesaria de palabras o construcciones. Sin
embargo, la normativa sigue una doble vara:

\begin{cita}
Si nos fijamos ahora en la aproximación retórica y gramatical a la redundancia,
observaremos una extraña dicotomía: las caracterizaciones de la redundancia que
suelen hacerse desde la filología dan a entender que se trata de una opción
legítima si procede de la pluma de los buenos escritores, a la vez que
ilegítima en los demás casos. La definición de \textit{pleonasmo} en el
\textit{DLE} ilustra esta dicotomía con claridad: «Empleo en la oración de uno o
más vocablos innecesarios para que tenga sentido completo, pero con los cuales
se añade expresividad a lo dicho, como en \textit{en fuga irrevocable huye la
hora}» \parencite[30]{bosque2020}.
\end{cita}

Según lo planteado, entonces, las repeticiones deben evitarse cuando reiteran la
información ya expresada en la misma oración \exref{ex:02-nuevamente}, pero son
aceptadas cuando refuerzan una idea o señalan diferencias entre cosas
semejantes \exref{ex:02-puertas}:

\begin{exe}
  \ex\label{ex:02-nuevamente} Así, si el acuífero cumple con las condiciones
      adecuadas, el agua puede \forma{volverse nuevamente} pura («¿Qué hacer
      ante la contaminación de aguas subterráneas?», Sondagua,
      \url{https://www.sondagua.cl/blog/que-hacer-ante-la-contaminacion-de-aguas-subterraneas/}).
  \ex\label{ex:02-puertas} Pasamos \forma{puertas y puertas y cuartos}, hasta que
      llegamos a una gran pieza oscura (A.~Uslar Pietri, «Maichak, el hombre del
      río», p.~170).
\end{exe}

En el caso de \exref{ex:02-matamos}, una estructura condicional tautológica
(\textit{NGLE}, §\,47.5e) pone de relieve la escasa importancia otorgada a lo que
le puede suceder (como si se dijera: \forma{si nos matamos, no pasa nada}) al
mismo tiempo que manifiesta que se trata de una situación irremediable.

\begin{exe}
  \ex\label{ex:02-matamos} Si nos matamos, nos matamos.
  \ex\label{ex:02-nopasa} \textit{No importa, no pasa nada}.
\end{exe}

En los ejemplos \exrefs{ex:02-carroza}{ex:02-nopasa} no se advierten problemas
respecto de la constitución interna de las oraciones: ninguna de las
consideraciones previas afecta la buena formación oracional.

\begin{ejercicio}\label{ej:02-elegancia}
Considere las siguientes oraciones y conteste las dos preguntas que se formulan
más abajo, como se muestra en el ejemplo resuelto~1.

\begin{exe}
  \ex No lo voy a repetir de nuevo.
  \ex Yo y mi hermano fuimos al cine.
  \ex Y si se comen unas patatas, se comen, y si no, no se comen (C.~Martín
      Gaite, \textit{El balneario}).
\end{exe}

\begin{consignas}
  \item ¿En qué se parecen y en qué contrastan estas tres oraciones? [La idea es
        considerar los elementos presentes o su disposición en la estructura.]
  \item ¿Qué problema plantea la estructura gramatical? ¿Se podría resolver de
        alguna manera lo observado en (a)?
\end{consignas}

\begin{resuelto}
1. (a) La oración está bien formada, aunque pueda señalarse como redundante (el
verbo ya encierra el valor de \sig{nuevamente}). (b) La oración no presenta
problemas gramaticales. La falta de elegancia se podría resolver ofreciendo
cambios léxicos o estructurales (\forma{No lo voy a repetir, No lo voy a volver
a decir, No lo voy a decir de nuevo}).
\end{resuelto}
\end{ejercicio}

\subsection{Gramaticalidad y significado referencial}

La buena formación de las expresiones de las lenguas naturales no depende de su
mayor o menor grado de relación con el referente, es decir, con aquella entidad
de la realidad extralingüística a la que remite:

\begin{cita}
Las secuencias que construimos combinando palabras pueden ser gramaticales y a
la vez relativamente incomprensibles, si su significación no se ajusta a lo que
independientemente sabemos sobre el mundo que nos rodea. Lo gramatical no es lo
\textsc{cierto}, ni siquiera lo \textsc{verosímil} [\ldots] Las secuencias que
se ajustan a los principios fundamentales de la gramática no tienen que denotar
situaciones ciertas o verosímiles \parencite[36]{bosque2009}.
\end{cita}

El contraste entre la gramaticalidad y la inconsistencia semántica fue
subrayado por Chomsky, al advertir que una frase como \forma{Las verdes ideas
incoloras duermen furiosamente} constituye, al mismo tiempo, una oración bien
formada desde el punto de vista sintáctico y un sinsentido desde la perspectiva
del significado. Sin embargo, lejos de representar una anomalía, esta clase de
contraste remite al \termino{desplazamiento}, una de las propiedades
constitutivas del lenguaje humano \parencite{hockett1971}, según la cual las
expresiones lingüísticas permiten hacer referencia tanto a lo que se encuentra
en el entorno inmediato (\forma{Hay un olor desagradable en la oficina}) como a
realidades distantes (\forma{Napoleón murió en Santa Elena en 1821}) o
inventadas (\forma{Superman nació en el planeta Krypton}). Así, el amplio
espectro de la ficcionalización hace interpretables oraciones que construyen no
solo sentidos paradójicos sino también mundos irreales:

\begin{cita}
Cuidamos las mancuspias hasta bastante tarde, ahora con el calor del verano se
llenan de caprichos y versatilidades, las más atrasadas reclaman alimentación
especial y les llevamos avena malteada en grandes fuentes de loza; las mayores
están mudando el pelaje del lomo, de manera que es preciso ponerlas aparte,
atarles una manta de abrigo y cuidar que no se junten de noche con las
mancuspias que duermen en jaulas y reciben alimento cada ocho horas
(J.~Cortázar, «Cefalea», p.~43).
\end{cita}

La contradicción (más en general, la oposición de sentidos, como en \forma{Se
produjo un silencio ensordecedor}), la autoconsistencia (es decir, la identidad
analítica de sentidos, como en \forma{Los triángulos tienen tres ángulos}), la
tautología (es decir, la unicidad en la asignación del valor de verdad, como en
\forma{Un peso es un peso}) y otros conceptos semejantes no entran en conflicto
con la noción de gramaticalidad, más bien constituyen «anomalías» desde el punto
de vista interpretativo. Por eso, la inserción de tales oraciones en contextos
situacionales muy específicos o en contextos lingüísticos ampliatorios puede
mejorar su significatividad (inferencia de implícitos, contrastes semánticos).
Así, la contradicción de \exref{ex:02-solteros-a} entre el significado de
\forma{soltero} (que no está casado) y \forma{casados} (que han contraído
matrimonio) se resuelve mediante el añadido de \forma{ahora}
\exref{ex:02-solteros-b}, haciendo referencia a dos momentos en el tiempo
(\forma{los que estaban solteros, ahora ya se han casado}). En
\exref{ex:02-catulo}, en el fragmento inicial del poema 85 de Catulo, la
contraposición entre los antónimos verbales \forma{odio} y \forma{amo} no
constituye en absoluto un sinsentido, sino que, por el contrario, refuerza lo
contradictorio del sentimiento amoroso.

\begin{exe}
  \ex\label{ex:02-solteros}
    \begin{xlist}
      \ex\label{ex:02-solteros-a} Todos los solteros de mi pueblo están casados.
      \ex\label{ex:02-solteros-b} Todos los solteros de mi pueblo están casados
          ahora.
    \end{xlist}
  \ex\label{ex:02-catulo} Odio y amo. ¿Por qué hago esto?, quizá preguntas. / No
      lo sé, pero siento que es así y me torturo.
\end{exe}

\begin{ejercicio}\label{ej:02-referencia}
Considere las siguientes oraciones y conteste las dos preguntas que se formulan
más abajo.

\begin{exe}
  \ex Una selenita nos atacó en la otra cuadra.
  \ex Asimismo agradecemos a Ireneo Fernando Cruz el habernos iniciado, durante
      un viaje a San Juan, en el conocimiento de las mancuspias (J.~Cortázar,
      «Cefalea», p.~43).
\end{exe}

\begin{consignas}
  \item ¿En qué se parecen y en qué contrastan estas dos oraciones?
  \item ¿Qué problema plantea la estructura gramatical? ¿Se podría resolver de
        alguna manera lo observado en (a)?
\end{consignas}
\end{ejercicio}

\begin{ejercicio}\label{ej:02-contradiccion}
Considere las siguientes oraciones y conteste las dos preguntas que se formulan
abajo, como se muestra en el ejemplo resuelto~1.

\begin{exe}
  \ex Los huérfanos de esta ciudad tienen unos padres maravillosos.
  \ex El ciego de la otra cuadra fue el primero que vio el accidente.
  \ex {[Un cliente llamando al mozo]} Este café caliente está frío.
  \ex El nuevo entrenador es simpático y antipático.
\end{exe}

\begin{consignas}
  \item ¿En qué se parecen y en qué contrastan estas oraciones? [Considere la
        manera en que se dan las relaciones de significados entre las palabras
        de las oraciones.]
  \item ¿Qué problema plantea la estructura gramatical? ¿Se podría resolver de
        alguna manera lo observado en (a)?
\end{consignas}

\begin{resuelto}
1. (a) La oración presenta estructuras bien formadas. Sin embargo, se establece
una contradicción entre lo que expresa la oración y el significado de una de las
palabras que contiene: \forma{huérfano} (persona menor de edad a quien se le han
muerto el padre y la madre o uno de los dos) es incompatible con el hecho de
\forma{tener unos padres maravillosos}. (b) No hay problemas de estructura
gramatical, aunque la contradicción es irresoluble porque entran en conflicto
los significados de las palabras: \forma{padres} y \forma{huérfanos}.
\end{resuelto}
\end{ejercicio}

\section{Gramaticalidad y variedades lingüísticas}

Finalmente, conviene observar que la gramaticalidad de las expresiones
lingüísticas, su buena formación, está referida a las distintas variedades de
una lengua, no se trata de una propiedad constitutiva abstracta. Si la gramática
y el léxico de una lengua tienen naturaleza mental (como se acepta desde
Saussure) e individual (como ha planteado Chomsky), es esperable que existan
diferencias explicables, tal vez, por el cambio en la disposición de las reglas
internas de la gramática (la modificación interna que se da en el dispositivo
gramatical). Bajo esta perspectiva, una lengua natural (el español, por
ejemplo) se debe entender como un conjunto de variedades lingüísticas o
lectos\footnote{El concepto de lecto se presentó en el \capref{1}.}
(cronolectos, dialectos, sociolectos e idiolectos, por usar denominaciones
clásicas), que comparten un altísimo número de principios y reglas
constitutivas internas y un número no menos importante de unidades léxicas
básicas. La gramaticalidad de las expresiones de una lengua, ante todo, tiene
que ver con sus distintas variedades: así, por ejemplo, ciertas variedades del
español aceptan la presencia del determinante delante del nombre propio
---\forma{la María, el Pedro}---, mientras que otras lo rechazan.

En términos generales, se pueden considerar cuatro tipos básicos de variaciones
lingüísticas: diacrónica, diatópica, diastrática y diafásica. Los conceptos
presentes en esta terminología, el primero proveniente de \textcite{saussure1945},
los tres restantes tomados de \textcite{coseriu1969}, remiten a las dimensiones
temporal, espacial, social y situacional en que se desarrollan los hechos del
lenguaje.

\subsection{Variación diacrónica}

Hace referencia a la dimensión temporal que origina las diferencias en el habla
entre las distintas generaciones, que se conocen como \termino{cronolectos}. Las
lenguas cambian a lo largo del tiempo, como muestran las dos versiones
siguientes del \textit{Auto de los Reyes Magos} (vv.~1-4 y 29-30, anónimo del
siglo~\textsc{xii}):

\begin{exe}
  \ex\label{ex:02-reyes}
    \begin{xlist}
      \ex\label{ex:02-reyes-a} ¡Dios criador, cuál maravila, / no sé cuál es
          aquesta strela! / Agora primas la e veída; / poco tiempo ha que es
          nacida. [\ldots] ¿En todo, en todo es nacido? / Non se si algo e veido
          (versión original).
      \ex\label{ex:02-reyes-b} ¡Dios Creador! ¡Qué maravilla! / ¿Qué estrella
          será esa que brilla? / Hasta ahora, no la he advertido; / hace bien
          poco que ha nacido. [\ldots] ¿Será verdad que ya ha nacido? / Dudo de
          lo que he advertido (versión de F.~Lázaro Carreter,
          1958).\pendiente{El original dice «L. Carreter»}
    \end{xlist}
\end{exe}

Además de las evidentes diferencias léxicas entre las dos versiones, ambas se
distinguen porque en \exref{ex:02-reyes-a}, que ejemplifica el castellano
medieval, los verbos con forma pasiva y significado activo, conocidos como
verbos deponentes en las gramáticas latinas (verbos inacusativos), tomaban como
auxiliar de los tiempos compuestos el verbo \forma{ser} (\forma{es nacida, es
nacido}), mientras que los demás verbos tomaban el auxiliar \forma{haber}
(\forma{e veída, e veido}). En el castellano moderno \exref{ex:02-reyes-b} esta
diferencia de auxiliares se perdió en favor del auxiliar \forma{haber}
(\forma{ha nacido, he advertido}).

\subsection{Variación diatópica}

Se trata de la variación que se produce en la dimensión geográfica entre
distintas zonas, regiones o lugares, y produce la diferencia entre
\termino{dialectos}, como las que se observan entre los ejemplos de
\exref{ex:02-queres} y \exref{ex:02-quetu}:

\begin{exe}
  \ex\label{ex:02-queres} ¿Qué querés vos para el 2021?
      (\url{https://www.instagram.com/reel/CFiNsjKltdb/}, del usuario comercial
      argentino @dalepaprika, de la red social Instagram).
  \ex\label{ex:02-quetu} ¿Qué tú quieres, mamita? ¿Qué tú quieres? («Qué tú
      quieres», del álbum \textit{Visualízate} de Gente de Zona, grupo cubano de
      salsa y reguetón).
\end{exe}

Mientras en la mayoría de las variedades del español el sujeto se pospone al
verbo en las oraciones interrogativas \exref{ex:02-queres}, en el español de
Cuba, el sujeto de las interrogativas mantiene la posición preverbal
\exref{ex:02-quetu}. Los ejemplos también muestran la diferencia entre el uso
del pronombre de segunda persona típico del español de Buenos Aires,
\forma{vos}, y la forma \forma{tú}, característica del español peninsular y de
la mayoría de las regiones hispanohablantes, como Cuba. Cada una de ellas
implica también diferencias en las formas verbales correspondientes:
\forma{vos querés / tú quieres}.

\subsection{Variación diastrática}

Constituye la dimensión referida a la estratificación o diferenciación social
referida a que las lenguas presentan variaciones relacionadas con los grupos o
estratos sociales, los \termino{sociolectos}, según la ocupación, profesión,
nivel de instrucción, entre otras variables sociales, tal como muestran las dos
interpretaciones siguientes del tango \textit{Cambalache} de Enrique Santos
Discépolo (1934):

\begin{exe}
  \ex\label{ex:02-sosa} ¡Que allá en el Horno / se vamo' a encontrar\ldots!
      (interpretación de J.~Sosa).
  \ex\label{ex:02-gardel} ¡Que allá en el Horno / nos vamo' a encontrar\ldots!
      (interpretación de C.~Gardel).
\end{exe}

En \exref{ex:02-sosa} y \exref{ex:02-gardel} se advierten diferencias en el
empleo del pronombre clítico\footnote{Los clíticos son pronombres personales
átonos que se apoyan para su pronunciación en la palabra que los precede o
sigue. Sus características y funciones se tratan en el \capref{10}.} de primera
persona plural. La forma \forma{se}, que aparece en \exref{ex:02-sosa}, estuvo
asociada durante décadas con los usos de los sectores populares urbanos,
mientras que \forma{nos}, la forma que aparece en \exref{ex:02-gardel}, siempre
se ha considerado como propia de la población escolarizada.

\subsection{Variación diafásica}\label{sec:02-diafasica}

Refiere a la dimensión vinculada con la situación social en que se usa la
lengua. Las formas lingüísticas se van adaptando a los distintos contextos
situacionales en que se produce el acto de habla, a las características del
entorno y del interlocutor, al tipo de discurso empleado y al grado de
formalidad de la comunicación. Tales aspectos constituyen el
\termino{registro}. Así, en \exref{ex:02-solicitud} y \exref{ex:02-guita} se
muestran dos situaciones diferentes en las que la misma persona, un alumno de la
carrera de Ingeniería Eléctrica de la Universidad Nacional del Sur, podría
formular un pedido:

\begin{exe}
  \ex\label{ex:02-solicitud} Por la presente, y en mi carácter de alumno de la
      carrera de Ingeniería Eléctrica, plan 2010, solicito una excepción que me
      permita cursar la asignatura que se detalla a continuación, sin cumplir
      con la totalidad de requisitos exigidos por el plan de estudios
      correspondiente (Solicitud de excepción para cursar o rendir materias,
      según el modelo de la UNS).
  \ex\label{ex:02-guita} Che, tengo que pedirte un favor muy grande. Este mes se
      me está haciendo muy largo y me estoy quedando corto de guita. ¿No me
      aguantás unos mangos hasta que cobre? En una semana te devuelvo todo.
\end{exe}

\begin{ejercicio}\label{ej:02-torito}
Lea atentamente el siguiente fragmento adaptado de \textit{Torito} de
J.~Cortázar y luego resuelva lo pedido más abajo, como se muestra en el ejemplo
resuelto.

\begin{cita}
Qué le vas a hacer, \subrayado{ñato}, cuando estás abajo todos te
\subrayado{fajan}. Todos, che, hasta el más maula. Te sacuden contra las sogas,
te encajan la \subrayado{biaba}. Vos te crees que yo me desespero, lo que pasa
es que no doy más aquí tumbado todo el día. Las noches de invierno son más
largas que esperanza'e pobre. Estoy todo el tiempo mirando el techo. Ahí tenés
otra cosa que no sé hacer, mirar p'arriba. Todos dijeron que me hubiera
convenido, que hice la gran \subrayado{macana} de levantarme a los dos
segundos, \subrayado{cabrero} como la gran flauta. Tienen razón, si me quedo
hasta las ocho no me agarra tan mal el rubio.
\end{cita}

\begin{consignas}
  \item Busque en el \textit{Diccionario de la lengua española} (\textit{DLE})
        de la Real Academia el significado de las palabras subrayadas y
        transcriba la acepción correspondiente a su uso en el texto.
  \item Una de las palabras subrayadas no está en este diccionario. ¿Por qué
        cree que sucede y en qué diccionario podría encontrarla?
  \item ¿Qué tipo de variación ejemplifican estas palabras?
\end{consignas}

\begin{resuelto}
\forma{ñato}: Am. chato (‖ de nariz poco prominente).
\end{resuelto}
\end{ejercicio}

\begin{ejercicio}\label{ej:02-variacion}
Considere las siguientes oraciones y resuelva lo pedido más abajo, como se
muestra en el ejemplo resuelto~1.

\begin{exe}
  \ex Hay que estudiar, \subrayado{chaval}, que la vida está muy mala
      (I.~Aldecoa).
  \ex ¿Es \subrayado{guapo}? No me he fijado en él.
  \ex Palermo, en estos días, era un \subrayado{rancherío} de muy mala muerte
      (M.~Caparrós, \textit{Página 12}, 1993).
  \ex Esto que \subrayado{agora} emprenderéis ---agregó el
      \subrayado{lectoral}--- será en servicio de la santa Iglesia de Cristo
      (E.~Larreta, \textit{La gloria de Don Ramiro}).
  \ex Ante el alza de los alquileres, algunos sectores del gobierno sugieren
      como alternativa el sistema de \subrayado{\textit{leasing}}.
  \ex Es una de las regiones más calurosas y \subrayado{paupérrimas} de todo el
      territorio nepalés.
\end{exe}

\begin{consignas}
  \item Busque en el diccionario el significado de las palabras subrayadas y
        transcriba la acepción correspondiente a su uso en las oraciones.
  \item Indique qué tipo de variación ejemplifica cada una de las palabras
        subrayadas.
\end{consignas}

\begin{resuelto}
1. (a) \forma{Chaval}: m. y f. coloq. Esp. Niño o joven. (b) Variación
dialectal.
\end{resuelto}
\end{ejercicio}

\begin{recuadro}[Reflexión final]
Aunque pueda resultar obvio, debería quedar claro que las tipologías anteriores
referidas a las diferentes variedades lingüísticas coexisten en cada individuo
que articula una oración en su variedad de lengua, una variedad que es a la vez
temporal (hablantes de diferentes generaciones), geográfica (hablantes de
diferentes regiones) y social (hablante con diferentes grados de instrucción o
con distintas ocupaciones), y que se produce en un determinado contexto
situacional (diferentes grados de formalidad y cuidado en el habla). Sin
embargo, muchas veces se pone el eje explicativo del lenguaje en la variedad por
la variedad misma, como si una sola dimensión determinara la conducta
lingüística de los hablantes y oyentes. Alternativamente, son los conocimientos
tácitos e interiores al individuo, su competencia lingüística (gramatical) y su
competencia comunicativa (situacional), los elementos que hacen posible una
conducta lingüística concreta por parte del hablante/oyente. En este sentido,
la competencia del individuo, protagonista central de la lengua, puede parecer
una explicación general y abstracta, pero resulta más esclarecedora que la mera
clasificación de variedades y lectos, que constituyen generalizaciones
exteriores y ajenas a los individuos.
\end{recuadro}

% --- Secciones finales -------------------------------------------------
\lecturasrecomendadas

\solucionario

\begin{solucion}{\ref{ej:02-definiciones}}
\begin{tabularx}{\linewidth}{@{}L *{4}{Y{.8cm}}@{}}
  \toprule
  De acuerdo con las definiciones dadas, la gramática de una lengua\ldots & 2 & 3 & 4 & 5 \\
  \midrule
  \ldots\ representa la descripción de un sistema de valores lingüísticos.    & \textbf{\menos} & \textbf{\mas} & \textbf{\mas} & \textbf{\menos} \\
  \ldots\ constituye un mecanismo capaz de producir estructuras lingüísticas. & \textbf{\menos} & \textbf{\menos} & \textbf{\menos} & \textbf{\mas} \\
  \ldots\ implica una regulación, exterior al individuo que la usa.           & \textbf{\mas} & \textbf{\menos} & \textbf{\menos} & \textbf{\menos} \\
  \ldots\ supone un sistema o mecanismo interior al individuo.                & \textbf{\menos} & \textbf{\mas} & \textbf{\mas} & \textbf{\mas} \\
  \ldots\ conforma una disciplina normativa o prescriptiva.                   & \textbf{\mas} & \textbf{\menos} & \textbf{\menos} & \textbf{\menos} \\
  \ldots\ conforma una disciplina descriptiva o científica.                   & \textbf{\menos} & \textbf{\mas} & \textbf{\mas} & \textbf{\mas} \\
  \bottomrule
\end{tabularx}
\pendiente{Falta la fila «remite al uso lingüístico de ciertos grupos sociales»}
\end{solucion}

\begin{solucion}{\ref{ej:02-searle}}
\begin{tabularx}{\linewidth}{@{}L Y{1.7cm} Y{1.7cm}@{}}
  \toprule
  ¿Qué tipo de regla encierran los siguientes enunciados? & Regulativa & Constitutiva \\
  \midrule
  b) No corresponde decir «pienso de que mañana empiezo la Facultad». & \textbf{\mas} & \textbf{\menos} \\
  c) No se dice «haiga», se dice «haya». & \textbf{\mas} & \textbf{\menos} \\
  d) «Si seguís hablando, la vas a pasar mal» se entiende como amenaza. & \textbf{\menos} & \textbf{\mas} \\
  e) Toda estructura lingüística contiene un núcleo. & \textbf{\menos} & \textbf{\mas} \\
  f) Usar el léxico lunfardo queda prohibido en la radiodifusión argentina. & \textbf{\mas} & \textbf{\menos} \\
  \midrule
  ¿Qué tipo de regla ofrecen las gramáticas normativas o prescriptivas? & \textbf{\mas} & \textbf{\menos} \\
  ¿Qué tipo de regla presentan las gramáticas descriptivas o científicas? & \textbf{\menos} & \textbf{\mas} \\
  ¿Qué tipo de regla es responsable de la regulación interna de una lengua? & \textbf{\menos} & \textbf{\mas} \\
  ¿Qué tipo de regla es responsable de la regulación externa de una lengua? & \textbf{\mas} & \textbf{\menos} \\
  \bottomrule
\end{tabularx}
\end{solucion}

\begin{solucion}{\ref{ej:02-correccion}}
\begin{tabularx}{\linewidth}{@{}L Y{1.5cm} Y{1.5cm}@{}}
  \toprule
  ¿Qué evaluación merecen las siguientes secuencias lingüísticas? & Correcta & Gramatical \\
  \midrule
  2. a. Las diputadas están saliendo y están entrando a otra reunión. & \textbf{\menos} & \textbf{\mas} \\
  \phantom{2.} b. Las diputadas están saliendo de una reunión y están entrando a otra. & \textbf{\mas} & \textbf{\mas} \\
  \addlinespace
  3. a. Trajimos tres remeras iguales para que no haya diferencias de color. & \textbf{\mas} & \textbf{\mas} \\
  \phantom{3.} b. Trajimos tres remeras iguales para que no haiga diferencias de color. & \textbf{\menos} & \textbf{\mas} \\
  \phantom{3.} c. Trajimos tres remeras iguales para que no hay diferencias de color. & \textbf{\menos} & \textbf{\menos} \\
  \addlinespace
  4. a. Tu mujer trabaja todo el día, parece \textit{blecansain}. & \textbf{\menos} & \textbf{\menos} \\
  \phantom{4.} b. Tu \textit{jermu} trabaja todo el día, parece incansable. & \textbf{\masmenos} & \textbf{\mas} \\
  \bottomrule
\end{tabularx}

2. a. [[Las diputadas] [están saliendo Ø] y [están entrando [a otra
reunión]].] La falta de la preposición locativa direccional \forma{de} no impide
la gramaticalidad de la oración, ya que su omisión se considera una cuestión
estilística. b. [[Las diputadas] [están saliendo [de una reunión]] y [están
entrando [a otra]]]. En el contraste de 2a-2b se observa que pueden existir
elementos no manifiestos (presupuestos) dentro de una estructura oracional bien
formada. Los verbos de movimiento como \forma{salir} y \forma{entrar} pueden
tener implícito el argumento locativo direccional que requieren porque se pueden
inferir del significado del mismo verbo, ya que \forma{salir} es \sig{egresar
del lugar donde se está} (presuposición de origen) y \forma{entrar} es
\sig{ingresar al lugar donde no se está} (presuposición de destino). A
diferencia del caso planteado en el ejemplo resuelto 1, en el punto b, en el que
la omisión de los argumentos del verbo produce una estructura agramatical, la
omisión de 2b no resulta en una estructura mal formada.

3. a. [Trajimos tres remeras iguales [para que no [haya] diferencias de
color].] b. [Trajimos tres remeras iguales [para que no [haiga] diferencias de
color].] c. [Trajimos tres remeras iguales [para que no [hay] diferencias de
color].] En el contraste de 3a-3b-3c se muestra el carácter histórico de la
normativa, en el sentido de que a lo largo del tiempo puede ir cambiando: en
3b, una forma existente en otras épocas como \forma{haiga}, hoy día resulta
incorrecta; en cambio, en 3a se evidencia la exigencia del modo subjuntivo en
construcciones finales, lo que torna agramatical la secuencia de 3c.

4. a. [Tu mujer trabaja todo el día, parece \juicio{*}{\forma{ble-cansa-in}}.]
b. [Tu \forma{jer-mu} trabaja todo el día, parece incansable.] En las formas de
4a-4b se extiende el concepto de gramaticalidad, inicialmente propio de la
sintaxis, a la morfología. En 4a, la estructura interna de la palabra
\forma{in-cansa-ble} aparece modificada, con las posiciones del prefijo y el
sufijo invertidas. En 4b la modificación afecta a la estructura silábica:
\forma{mu-jer} se invierte y el resultado es una combinación apropiada
\forma{jer-mu}, frente a otras que serían imposibles para el español, como
\juicio{*}{\forma{jre-um}}, \juicio{*}{\forma{um-jre}}.

\textit{Preguntas sobre 4a y 4b}. (4a) La forma
\juicio{*}{[\forma{ble-cansa-in}]}, además de que carece de un significado
predecible, viola las reglas de formación de palabras (los prefijos deben ir
delante de la raíz y los sufijos, detrás). Esto es el opuesto exacto de lo que
ocurre con la forma \forma{in-cansa-ble}, no solo calculable en cuanto a su
significado, sino también porque responde a las reglas de formación de palabras
en español. (4b) En contraste con lo anterior, ambas formas \forma{mu-jer} y
\forma{jer-mu} tienen viabilidad fonológica, pero difieren en el registro
(v.~§\ref{sec:02-diafasica}): la primera corresponde al español general y la segunda es una
palabra propia del lunfardo (léxico popular argentino). En \forma{jermu} se
aplica el procedimiento típico del lunfardo conocido como \termino{vesre} o
inversión silábica, como en \forma{feca} (\forma{café}), \forma{gotán}
(\forma{tango}), \forma{gomía} (\forma{amigo}). Por otra parte, la forma
\forma{jermu} tiene una interpretación más específica o restringida, ya que no
refiere a cualquier \forma{mujer}, sino a la «esposa o pareja de alguien».
\end{solucion}

\begin{solucion}{\ref{ej:02-registros}}
\begin{consignas}
  \item En 4a, la forma \juicio{*}{[\forma{ble-cansa-in}]} resulta anómala para
        cualquier hablante del castellano, frente a la bien formada
        [\forma{in-cansa-ble}]. En 4b, la forma [\forma{jer.mu}], dado que
        corresponde al léxico lunfardo, tiene aceptación restringida al registro
        coloquial del Río de la Plata.\pendiente{La consigna pregunta por 3a y 3b}
  \item Condiciones de uso, registros y estilos de lengua remiten al ámbito de
        la corrección, ya que lo agramatical, dado que contraviene las reglas
        internas o constitutivas del sistema, es lo que «en condiciones
        normales» el hablante no produce. Así, la gramaticalidad de
        [\forma{jer.mu}], cuyo empleo está fehacientemente atestiguado, se
        encuentra salvada frente a la agramaticalidad de
        \juicio{*}{[\forma{ble-in-cansa}]}, que no sería una forma que los
        hablantes emitirían naturalmente.
\end{consignas}
\end{solucion}

\begin{solucion}{\ref{ej:02-aceptabilidad}}
\begin{tabularx}{\linewidth}{@{}L Y{1.5cm} Y{1.5cm}@{}}
  \toprule
  ¿Qué evaluación merecen las siguientes secuencias lingüísticas? & Gramatical & Aceptable \\
  \midrule
  2. La madre del primo de la novia del hermano de mi suegra está contenta. & \textbf{\mas} & \textbf{\masmenos} \\
  \addlinespace
  3. a. «Toda fuerza armada o reunión de personas que se atribuya los derechos
  del pueblo y peticione a nombre de éste, comete delito de sedición» (CN,
  art.~22). & \textbf{\mas} & \textbf{\mas} \\
  \phantom{3.} b. «Toda fuerza armada o reunión de chabones que se atribuya los
  derechos del pueblo y peticione a nombre de éste, comete delito de sedición»
  (¿CN, art.~22?). & \textbf{\mas} & \textbf{\menos} \\
  \bottomrule
\end{tabularx}

2. La estructura de 2 (como en el ejemplo resuelto 1) no presenta problemas de
gramaticalidad, pero la aceptabilidad resulta al menos dudosa: en cada instancia
se introducen sustantivos que remiten a nuevos referentes (\forma{madre, primo,
novia, hermano, suegra}), cuya acumulación dificulta el procesamiento cognitivo
de la oración. [[La madre [del primo [de la novia [del hermano [de mi
suegra]]]]] está contenta.]

3. La selección léxica de los textos formales, como el de la Constitución
Nacional, explica las diferencias en la aceptabilidad de 3a frente a 3b. En
ambos casos la buena formación oracional está asegurada, solo que el lexema
\forma{chabones} en 3b la torna inaceptable por cuestiones de estilo;
\forma{chabones} es una palabra demasiado coloquial para el registro que
corresponde a un texto formal como la Constitución. En contraste, la forma
\forma{personas} en 3a es la esperable dentro de un texto de naturaleza
jurídica.
\end{solucion}

\begin{solucion}{\ref{ej:02-ingrid}}
\begin{consignas}
  \item La estructura de las oraciones, en todas las repeticiones que se entonan
        en la canción, resulta gramatical; no se advierten formas incompletas,
        como podría ser la falta de algún argumento verbal. Lo que se advierte
        son problemas en la construcción del discurso, lo que se conoce como
        anacoluto, y también arranques de frases que no se completan y que se
        retoman desde el inicio con una nueva oración.
  \item La aceptabilidad de las oraciones formadas resulta dispar debido a la
        selección léxica algo errática que se produce en las últimas dos
        repeticiones de la canción. Con todo, se observa que la elección de
        palabras, aunque presenta variaciones, se produce dentro del mismo campo
        semántico de la forma original; así, según se puede escuchar, aparece
        \forma{ventilador} en lugar de \forma{aire (acondicionado)} y
        \forma{cerrar}, en vez de \forma{abrir (la botella)}.
\end{consignas}
\end{solucion}

\begin{solucion}{\ref{ej:02-ambiguedad}}
(a) Las tres oraciones del ejercicio están bien formadas, son gramaticales,
aunque ambiguas, ya que, fuera de contexto, admiten múltiples significados. La
ambigüedad se relaciona con factores sintácticos en el caso resuelto 1 y en 3,
y con factores léxicos en 2.

2. a. \textit{Antes de salir a navegar, acomodaron correctamente las velas (que
utilizaban para alumbrar)}. b. \textit{Antes de salir, acomodaron correctamente
las velas (que servían para impulsar el barco)}.

3. a. [Yo [no recibiría] a tu hermano en mi oficina.] a'. \textit{Yo no
recibiría a tu hermano en mi oficina (es más, lo ignoraría)}. b. [[Yo] no
recibiría a tu hermano en mi oficina.] b'. \textit{Yo no recibiría a tu hermano
en mi oficina (pero sí mi secretario)}. c. [Yo no recibiría [a tu hermano] en mi
oficina.] c'. \textit{Yo no recibiría a tu hermano en mi oficina (pero sí a su
enviado)}. d. [Yo no recibiría a tu hermano [en mi oficina].] d'. \textit{Yo no
recibiría a tu hermano en mi oficina (pero sí en mi casa)}.

(b) Las estructuras no presentan problemas gramaticales, y la ambigüedad
interpretativa (multiplicidad de significados) deriva de la combinatoria
sintagmática o léxica. En 2, se podrían establecer alternativas léxicas o
especificaciones para desambiguar los lexemas: \forma{Antes de salir a navegar,
acomodaron correctamente las velas para alumbrar} frente a \forma{Antes de salir
a navegar, acomodaron correctamente el velamen del barco}. En 3 se podría dar
una marcación especial al constituyente focalizado,\footnote{El foco es el
segmento sintáctico del enunciado al que se otorga mayor relieve informativo
(\textit{DLE}, en línea). Este concepto se introduce en el \capref{6}.} tal
como ocurre en la oralidad, en la que se enfatiza con la entonación el
constituyente sobre el que se centra la interpretación, frente a la entonación
neutra que caracteriza a la negación del núcleo predicativo: \forma{Yo no
recibiría a tu hermano en mi oficina, \textup{YO} no recibiría a tu hermano en
mi oficina, Yo no recibiría \textup{A TU HERMANO} en mi oficina, Yo no
recibiría a tu hermano \textup{EN MI OFICINA}}.
\end{solucion}

\begin{solucion}{\ref{ej:02-elegancia}}
(a) Las oraciones 2 y 3 están bien formadas. 2 puede señalarse como descortés
(sobre este tipo de oraciones pesa aquello de «el burro adelante, para que no se
espante»). La elegancia de una expresión lingüística, en todo caso, no tiene que
ver con la gramática y la estructura de la lengua, sino con una valoración de
naturaleza social. En el caso de 3, lo que se quiere indicar con la repetición
es la importancia de tomar una decisión, sea esta la que fuera (\textit{NGLE},
§\,48.5e). (b) Las estructuras no presentan problemas gramaticales. La falta de
elegancia se podría resolver en 2 restableciendo el orden de los constituyentes
(\forma{Mi hermano y yo fuimos al cine, Fui al cine con mi hermano}). La
tautología de 3, por su parte, se resolvería con cambios en la estructura
(\forma{Y si se comen unas patatas, todos felices, y si no, también})
sacrificando el valor de expresividad que añade la repetición.
\end{solucion}

\begin{solucion}{\ref{ej:02-referencia}}
(a) Las oraciones 1 y 2 están bien formadas, aunque las referencias de
\forma{selenita} (habitante imaginario de la luna) y de \forma{mancuspia}
(palabra inventada por Cortázar) no remitan a entidades reconocidas en el mundo
real: la buena formación oracional no depende de su mayor o menor veracidad, es
decir, de que su contenido sea cierto o por lo menos verosímil («se non è vero,
è ben trovato»). Ni las ficciones de la vida cotidiana ni la ficción literaria
(la de Cortázar, en el ejemplo) serían posibles si la viabilidad de las
oraciones dependiera de su veracidad. (b) Las estructuras están bien formadas,
carecen de problemas gramaticales. La escasa verosimilitud de 1 y la falta de
certidumbre respecto de a qué remite \forma{mancuspias} en 2 se podrían
resolver introduciendo otros lexemas que, aunque cambien o precisen el
significado, permitirían corroborar que ambas oraciones están bien formadas. Así
en 1 se podría reemplazar \forma{selenita} por \forma{delincuente}: \forma{Un
delincuente nos atacó en la otra cuadra}; y en 2, en vez de \forma{mancuspias},
se podría colocar \forma{araucarias}: \forma{Asimismo agradecemos a Ireneo
Fernando Cruz el habernos iniciado, durante un viaje a San Juan, en el
conocimiento de las araucarias}.
\end{solucion}

\begin{solucion}{\ref{ej:02-contradiccion}}
(a) Las oraciones 2-4 presentan estructuras bien formadas. Sin embargo, en 2 (al
igual que en el ejemplo resuelto 1) se establecen contradicciones entre lo que
expresa la oración y el significado de una de las palabras que contienen. En 2,
la incompatibilidad se produce entre \forma{ciego} (privado de la vista) y el
hecho de \forma{ver el accidente}. En 3 se produce una contraposición entre
antónimos, en este caso, \forma{caliente} y \forma{frío}. En 4 la
contradicción se da entre una palabra y su antónimo. (b) No hay problemas de
estructura gramatical en los ejemplos. En 2 la contradicción es irresoluble
porque entran en conflicto los significados de las palabras: \forma{ciego} y
\forma{vio}. Los casos de 3 y 4 podrían ser salvables. En 3 el contexto
situacional permitiría interpretar la contradicción como aparente asignando un
sentido irónico a las palabras del cliente, que se estaría quejando porque el
café caliente que pidió, el mozo se lo trajo frío. En 4 la contradicción
también se puede resolver interpretando que el nuevo entrenador en ocasiones es
simpático y en otras no, señalando que se trata de una característica variable
de la persona de la que se habla, que puede depender de las circunstancias o del
momento.
\end{solucion}

\begin{solucion}{\ref{ej:02-torito}}
(a) \forma{Fajar}: tr. Can. y Am. Pegar a alguien, golpearlo. \forma{Biaba}: no
se encuentra en el \textit{DLE}. \forma{Macana}: f. coloq. Arg., Bol., Par.,
Perú y Ur. Mentira, desatino. \forma{Cabrero}: adj. coloq. Arg. Que es, o suele
estar, malhumorado, agresivo o irritado. (b) La palabra que no está en el
\textit{DLE} es \forma{biaba}. Se puede encontrar en un diccionario del lunfardo
como el \textit{DEL} \parencite{conde2004} y significa \sig{golpe}. (c)
\forma{Chato} y \forma{fajar} son americanismos, \forma{macana} se usa en varios
países de Sudamérica y \forma{cabrero} es un argentinismo; por lo tanto,
ejemplifican la variedad dialectal. El lunfardismo \forma{biaba} ejemplifica la
variedad diastrática.
\end{solucion}

\begin{solucion}{\ref{ej:02-variacion}}
(a) 2. \forma{Guapo}: adj. Bien parecido. 3. \forma{Rancherío}: m. Arg., Chile
y Ur. Conjunto desordenado y miserable de viviendas precarias en las zonas
rurales. 4. \forma{Agora}: adv. dem. desus. ahora (‖ en este momento). U. c.
rur. \forma{Lectoral}: canónigo lectoral. m. Prebendado de oficio. Es el
teólogo del cabildo, y deberá ser licenciado o doctor en teología. 5.
\forma{Leasing}: Econ. Arrendamiento con opción de compra del objeto arrendado.
6. \forma{Paupérrimo}: adj. sup. cult. de \forma{pobre}. (b) 2 y 3: variación
dialectal. 4: variación diacrónica (\forma{agora}). 5 y 6: variación
diastrática.
\end{solucion}

